\uuid{7fpW}
\titre{Intégration par parties}
\auteur{hanani ; mzari}
\organisation{exo7}
\licence{CC BY-NC-SA 4.0}
\datecreate{2020-01-20}
\niveau{L1}
\module{Analyse}
\chapitre{Calcul d'intégrales}
\sousChapitre{Autre}
\difficulte{1}
\theme{Calculs d'intégrales, intégration par parties}
\exo7qcmid{lille-2.integrales.003}

\contenu{
\begin{question}
\qclassification{127.99}
\qtags{motcle=intégration par parties}
Parmi les égalités suivantes, cocher celles qui sont vraies :
\begin{answers}
\bad{$\displaystyle \int _0^{\pi}x\sin x\,\mathrm{d}x=\int _0^{\pi}\cos x\mathrm{d}x$.}
\good{$\displaystyle \int _0^{\pi}x\sin x\,\mathrm{d}x=\pi +\int _0^{\pi}\cos x\mathrm{d}x$.}
\bad{$\displaystyle \int _0^{\pi}x\sin x\,\mathrm{d}x=\pi -2$.}
\good{$\displaystyle \int _0^{\pi}x\sin x\,\mathrm{d}x=\pi$.}
\end{answers}
\vskip2mm
\begin{explanations}
Une intégration par parties, avec $u=x$ et $v=-\cos x\Rightarrow v'=\sin x$, donne
$$\int _0^{\pi}x\sin x\,\mathrm{d}x=\Big[-x\cos x\Big]_0^{\pi}+\int _0^{\pi}\cos x\mathrm{d}x=\pi +\int _0^{\pi}\cos x\mathrm{d}x=\pi.$$
\end{explanations}
\end{question}
}
